\section{The Capacity-Radius Regime}

\subsection{The finite Hamming grid}

Let \(1\le k\le n\).  The capacity endpoint is already an attainable
normalized Hamming radius:
\[
  r_{\mathrm{cap}}=n-k,
  \qquad
  \deltacap=\frac{n-k}{n}=1-\rho.
\]
There is therefore no rounding loss at the endpoint.  For example, the
rate-\(1/4\), length-\(4096\) member has
\[
  k=1024,\qquad r_{\mathrm{cap}}=3072,\qquad
  \deltacap=\frac34.
\]
The endpoint theorem below includes the strictly smaller radius
\(3071/4096\).

There is, however, a distinct granularity issue in the phrase ``largest
relative radius.''  A closed Hamming ball changes only at
\[
  \mathcal H_n=\left\{0,\frac1n,\ldots,\frac nn\right\}.
\]
Consequently, a safe-radius set defined over all real
\(\delta\in[0,1]\) may be half-open and have no largest element.  We report
both the largest radius in \(\mathcal H_n\) and the supremum over real radii.

\subsection{An explicit smooth large-field family}
\label{sec:explicit-family}

For \(u\in\mathbb N\), set
\[
  k_u=64\cdot2^u=2^{u+6},
  \qquad
  n_u=4k_u=256\cdot2^u=2^{u+8},
  \qquad
  q_u=3^{n_u}.
\]
Let \(\F_u=\mathbb F_{q_u}\).  Since the multiplicative group of a finite
field is cyclic of order \(q_u-1\), it suffices to check
\(n_u\mid(q_u-1)\).  Euler's theorem gives this directly:
\[
  \varphi(n_u)=\frac{n_u}{2},
  \qquad
  \gcd(3,n_u)=1,
  \qquad
  3^{n_u/2}\equiv1\pmod{n_u}.
\]
Squaring yields \(3^{n_u}\equiv1\pmod{n_u}\).  Thus
\(\F_u^\times\) contains a subgroup \(H_u\) of order \(n_u\).  We take
\(\mathcal L_u=H_u\) and
\[
  \mathcal C_u=\RS[\F_u,H_u,k_u].
\]
The evaluation length and degree bound are powers of two, the domain is a
multiplicative subgroup, and the rate is exactly \(1/4\).  This is the
standard smooth-domain condition for FFT-based Reed--Solomon codes
\cite{arnonChiesaFenziYogev2024}, and it is the explicit family used in the
quantitative statement.

\subsection{The field-size inequality}
\label{sec:field-inequality}

Fix \(\epsstar=2^{-128}\).  For the family above,
\[
  \binom{n_u}{k_u}\le 2^{n_u}.
\]
Moreover \(n_u\) is even and \(n_u/2\ge128\), so
\[
\begin{aligned}
  q_u=3^{n_u}
  &=(3^2)^{n_u/2}
    =9^{n_u/2} \\
  &\ge8^{n_u/2}
    =2^{3n_u/2}
    \ge2^{n_u+128}.
\end{aligned}
\]
It follows that
\[
  \binom{n_u}{k_u}
  \le 2^{-128}q_u
  =\epsstar|\F_u|.
\]

The pair-MCA theorem uses the comparison
\[
  2^{128}\binom{n_u}{k_u+1}\le q_u.
\]
It follows from the same calculation, because
\(\binom{n_u}{k_u+1}\le2^{n_u}\).
Compared with the ordered certificate count, this removes a factor \(k_u+1\):
\[
  \binom{n_u}{k_u}(n_u-k_u)
  =(k_u+1)\binom{n_u}{k_u+1}.
\]
The field is exponential in the block length.  This is a valid explicit
sufficiently-large-field regime, not a claim about practically sized fields
or about every smooth code supplied to the Grand Challenge.
